% =====================================================================================
\chapter{METHODOLOGY: PERCEPTION-GUIDED ROBOTIC TACK PLACEMENT}
\label{chp:method_cell}
% First shape 2026-10-16 (from the 2026-10-08 skeleton chapters 4-6 + the work since).
% Shape follows the lab examples (Oydu, Kurt, Agus): overview -> problem formulation ->
% one section per component, its mathematics in place -> implementation details.
% Material: README (Big picture, How the pieces connect, Algorithms §1-8, §14-16, node
% table); notes/perception_aciklama.md, realtime_fp.md, multiview_refine_plan.md,
% curved_seams_plan.md, seam_two_modes_plan.md, pen_marking_plan.md, aps_transit_plan.md,
% thesis_notes.md; module docstrings (named in each \SRC).
% Results do NOT go here: every number measured on the bench belongs to Chapter 6.
% =====================================================================================

This chapter describes the robotic cell that turns a hand-placed assembly of known parts
into marked tack positions, using one wrist-mounted RGB-D camera and no touch sensing
for localisation. Section~\ref{sec:mc_overview} gives the pipeline as a whole and
Section~\ref{sec:mc_problem} states the problem. The sections after them follow the
pipeline: calibration (\S\ref{sec:mc_calibration}), the operator interface
(\S\ref{sec:mc_hri}), perception and registration (\S\ref{sec:mc_perception}),
multi-view refinement (\S\ref{sec:mc_multiview}), seams (\S\ref{sec:mc_seams}), tacks
(\S\ref{sec:mc_tacks}) and motion (\S\ref{sec:mc_motion}). Implementation details close
the chapter (\S\ref{sec:mc_impl}).

% -------------------------------------------------------------------------------------
\section{Overview of the Proposed System}
\label{sec:mc_overview}
The cell works on assemblies of parts whose CAD models are known: plates, tubes and
curved strips placed by hand on a table, in weld position, one after the other. It
finds where they are, computes where they meet, places tacks on the weldable seams and
marks each tack with a pen held by a UR5e. The pen stands in for a welding torch: the
geometry, the reach and the accuracy requirements are the same, the heat is not.

The pipeline has four stages (Figure~\ref{fig:mc_pipeline}):
\begin{enumerate}
\item \textbf{Perceive.} The operator aims the camera by pointing, freezes a frame and
  clicks the part with gestures; the part is segmented, classified against the CAD
  library and registered in 6D, then tracked while it is placed and saved when it rests.
\item \textbf{Refine.} With every part saved, the robot moves the camera to a few
  close-range views around the joints and refines all poses jointly.
\item \textbf{Compute.} Each registered part is represented as an exact geometric
  primitive; seams follow from where the primitives meet, and tacks from a placement
  rule on the seams.
\item \textbf{Act.} Reachability, collision-free transits and a force-gated descent
  bring the pen to every tack.
\end{enumerate}

\figplaceholder{fig:mc_pipeline}{The pipeline: aim, capture, segment + classify,
register, track, save (per part) -> multi-view refinement -> seams -> tacks ->
reachability -> transits -> force-gated marking; which machine runs what}{Overview of the
proposed perception-guided tack placement pipeline.}
\SRC{notes/pipeline\_map.dot (redraw in English); README ``Big picture''.}

% -------------------------------------------------------------------------------------
\section{Problem Formulation}
\label{sec:mc_problem}
\textbf{Given} a library of CAD models $\{\mathcal{M}_m\}$, a manipulator with joints
$\vect{q}\in\mathbb{R}^6$ and forward kinematics $\Tf{B}{E}(\vect{q})$, an RGB-D camera
on the wrist, and an assembly of parts $k=1,\dots,K$ placed by hand, \textbf{find} for
every part its model $m_k$ and pose $\Tf{B}{O_k}$; from them the set of weldable seams
$\{\mathcal{S}_j\}$; on them a set of tacks, each a point $\vect{p}_i$ on a seam with a
pen axis $\vect{a}_i$; and for every tack a collision-free joint trajectory that brings
the pen tip to $\vect{p}_i$ along $\vect{a}_i$.

The quantity that matters is the mark error: the distance between where the pen marks
and the true weld root. It is the sum of the errors of every link in the chain
\begin{equation}
\Tf{B}{O}(t) = \Tf{B}{E}\big(\vect{q}(t)\big)\;\Tf{E}{C}\;\Tf{C}{O}(t),
\label{eq:frame_chain}
\end{equation}
which is why calibration, registration and refinement each get a section below.

\textbf{Assumptions.} The parts are rigid and match their CAD models; they rest on the
table or on each other (they do not float); the table is a plane; the operator places
the parts one at a time; nothing moves while the robot marks.

\TODO{Notation table (frames B, E, C, O; $\vect{q}$; seam $\mathcal{S}$, tack
$(\vect{p},\vect{a},\varphi)$ with roll $\varphi$; tolerances). One table, used by every
chapter.}

% -------------------------------------------------------------------------------------
\section{Coordinate Frames and Calibration}
\label{sec:mc_calibration}
The principle of the calibration is to measure each link of Equation~\ref{eq:frame_chain}
against a reference that does not depend on it. The robot is that reference: its
factory-calibrated kinematics and the pen tip, touched onto surfaces, measure better
than any camera-based estimate. The order matters too, because a step calibrated on top
of an uncorrected earlier one learns that error.

\begin{table}[htbp]
\caption{The calibration chain: each link, its method and its reference.}
\label{tab:calib_chain}
\centering\small
\begin{tabular}{@{}lll@{}}
\toprule
Link & Method & Reference \\
\midrule
Kinematics & factory calibration used in every tool & the controller \\
Pen tip (TCP) & 4-point TCP; roll and tilt touches & table plane \\
Table plane & pen touches, plane fit & the robot \\
Camera depth & Tare with a robot-computed distance & table plane \\
Extrinsic rotation & camera tilt separated from table tilt at four wrist yaws & table plane \\
Extrinsic translation & online, from multi-view disagreement (\S\ref{sec:mv_selfcal}) & views \\
\bottomrule
\end{tabular}
\end{table}
\TODO{Residual column moves to Chapter 6 (it is a result).}

\figplaceholder{fig:frames}{TF tree base\_link -- tool0 -- camera\_color\_optical\_frame
-- object, drawn on the robot}{Coordinate frames of the cell.}

\subsection{Kinematics, Tool Centre Point and Table}
\label{sec:mc_kin_tcp_table}
\TODO{Why the factory-calibrated chain (nominal vs.\ calibrated differ by millimetres at
the tip); the pen TCP from four orientations; the table as a plane
$\vect{n}^{\top}\vect{x}+d=0$ fitted to pen touches, tilted (it is not level), used as the
height reference everywhere (ground cut, depth check, collision model).}
\SRC{README \S14; scripts/table\_touchoff.py, tcp\_offset\_probe.py; notebooks/table\_plane.json.}

\subsection{Hand-Eye Calibration}
\label{sec:mc_handeye}
The camera's pose on the flange, $\mathbf{X}=\Tf{E}{C}$, solves
\begin{equation}
\mathbf{A}_i\,\mathbf{X} = \mathbf{X}\,\mathbf{B}_i,
\label{eq:handeye}
\end{equation}
with $\mathbf{A}_i$ the relative flange motions and $\mathbf{B}_i$ the relative motions of
a calibration board seen by the camera \cite{tsai1989handeye,park1994handeye}.
\TODO{The ChArUco board, the sample poses, the residual judge; the TCP-frame pitfall
(a solve in the wrong frame that looked plausible); separating the camera's tilt from
the table's by viewing the table at four wrist yaws.}
\SRC{README \S14 (24 and 28 Sep); helper/calibration/; thesis\_notes.md.}

\subsection{Depth Calibration Against the Robot}
\label{sec:mc_depth}
A disparity offset $\delta$ (pixels) in a stereo camera with focal length $f$ and baseline
$b$ gives a depth error that grows with the square of the range:
\begin{equation}
\Delta z = \frac{z^{2}\,\delta}{f\,b}.
\label{eq:depth_error}
\end{equation}
The camera's own calibration (the Tare) needs a ground-truth distance; here it is
computed from the robot (the camera pose by Equation~\ref{eq:frame_chain} and the touched
table plane) instead of measured by hand.
\TODO{Measurement at three heights; the $a + k z^2$ fit; why a hand-measured distance
left a scale error.}
\SRC{README \S16; scripts/tare\_distance.py, extrinsic\_check.py.}
\figplaceholder{fig:depth_law}{Camera-vs-pen table offset at three ranges, before and after
the Tare, with the $a+kz^2$ fit}{Depth error against range, before and after calibration.}

% -------------------------------------------------------------------------------------
\section{Operator Interaction}
\label{sec:mc_hri}
The operator never touches a keyboard: hand gestures seen by the same RealSense drive
the pipeline, and pointing aims the camera before a frame is captured.

\subsection{Gesture Interface}
\label{sec:mc_gestures}
MediaPipe's gesture recogniser \cite{lugaresi2019mediapipe} classifies seven static
gestures per frame. A gesture acts only after it has been stable for a number of frames
(a latch), and a gesture that moves the robot must be held for two seconds. Because the
hand and the parts are in one image, the index fingertip is a pixel of the very frame
that is segmented: no second camera, no calibration between cameras.
\begin{table}[htbp]
\caption{Gestures and their actions.}
\label{tab:gestures}
\centering\small
\begin{tabular}{@{}p{0.24\textwidth}p{0.34\textwidth}p{0.34\textwidth}@{}}
\toprule
Gesture & No frame frozen & A frame frozen (segmenting) \\
\midrule
Pointing\_Up & aim the camera (\S\ref{sec:mc_aim}) & cursor: dwell to pick a point \\
ILoveYou & capture: freeze an RGB-D frame & re-capture \\
Thumb\_Up & --- & commit the point, segment and register \\
Thumb\_Down & clear the frame & mark the point ``not this part'' \\
Open\_Palm & compute the seams & undo the last point \\
Closed\_Fist (held 2 s) & multi-view refinement (the robot moves) & same \\
Victory & save the part & save the part \\
\bottomrule
\end{tabular}
\end{table}
\SRC{scripts/gesture\_control\_node.py docstring; README Data flow \S E.}

\subsection{Finger-Pointing Camera Aim}
\label{sec:mc_aim}
While the operator points and nothing is frozen, the camera turns until the fingertip is
in the middle of the image. The motion is restricted to a rotation about the camera's
own optical centre: the camera's position is held, so the operator steers only where it
looks. This restriction is what makes the controller simple, because the ray from the
optical centre to a still fingertip does not change in the base frame however the
camera turns, whatever the fingertip's depth.

\paragraph{Error with exact latency compensation.} The fingertip pixel
$(u,v)$, measured on an image taken at time $t_k$ and delivered about 0.1~s later, is
turned into a ray in the base frame with the camera orientation at $t_k$ (from the
joint-state history), filtered there by a One-Euro filter \cite{casiez2012oneeuro}, and
expressed in the camera frame of the current control cycle:
\begin{equation}
\vect{r}_B = \Rf{B}{C}(t_k)\,\frac{\mathbf{K}^{-1}[u\ v\ 1]^{\top}}{\norm{\mathbf{K}^{-1}[u\ v\ 1]^{\top}}},
\qquad
\vect{d} = \Rf{B}{C}(t)^{\top}\,\vect{r}_B,
\qquad
\vect{e} = \big(d_x/d_z,\ d_y/d_z\big).
\label{eq:aim_error}
\end{equation}

\paragraph{Control law.} An image-based visual servo \cite{chaumette2006vs} in its
simplest form, a proportional loop on the angular error, with a deadband and a rate cap:
\begin{equation}
\boldsymbol{\omega}_C = \big(-K e_y,\ K e_x,\ -k_r\,\psi\big),
\qquad
\psi = \operatorname{twist}_z\!\big(\Rf{B}{C_0}^{\top}\Rf{B}{C}\big),
\label{eq:aim_law}
\end{equation}
where $\psi$ is the roll of the camera about its optical axis relative to the home view,
taken from a swing--twist split. Two rotations aim the optical axis; the third, the
roll, is redundant and is held at home so that the image keeps its orientation.
No integral term (it winds up over the latency) and no derivative term (the landmark
is noisy).
\TODO{Deadband with hysteresis; the rate and acceleration limits and the two-lag
smoother; the excursion limits looked ahead by the stopping distance.}

\paragraph{From the camera rate to joint rates.} The tool flange moves so that the optical
centre stays still, with a small proportional pull back to its anchor $\vect{p}_{C,0}$;
damped least squares \cite{wampler1986dls} gives the joint rates:
\begin{equation}
\begin{aligned}
\vect{v}_E &= k_p\,(\vect{p}_{C,0}-\vect{p}_C) + \boldsymbol{\omega}\times(\vect{p}_E-\vect{p}_C),
\qquad \boldsymbol{\omega} = \Rf{B}{C}\,\boldsymbol{\omega}_C,\\
\dot{\vect{q}} &= \mathbf{J}^{\top}\big(\mathbf{J}\mathbf{J}^{\top}+\lambda^2\mathbf{I}\big)^{-1}
\begin{bmatrix}\vect{v}_E\\ \boldsymbol{\omega}\end{bmatrix}.
\end{aligned}
\label{eq:aim_dls}
\end{equation}

\paragraph{Execution and safety.} The joint rates stream to the UR driver's velocity
controller, which the aim switches in and out by itself. Every cycle, the pose 0.3~s
ahead is checked against the collision model, the branch lock and the wrist singularity;
a fingertip lost for 0.3~s ramps the rate to zero.

\paragraph{Command interpolation.} The control law runs at 125~Hz (one cycle costs
2.5~ms in Python); the controller runs at 500~Hz. A velocity held for several controller
cycles and then changed is a step that the arm applies as a torque kick; a stream of
them is audible. A relay between the two ramps linearly from one command to the next
over the measured command period $T$,
\begin{equation}
\dot{\vect{q}}_{\mathrm{cmd}}(t) = \dot{\vect{q}}_{k-1} + \frac{t-t_k}{T}\,\big(\dot{\vect{q}}_k-\dot{\vect{q}}_{k-1}\big),
\qquad t_k \le t < t_k+T,
\label{eq:relay}
\end{equation}
and stops the arm by itself when the commands stop.
\figplaceholder{fig:mc_aim_loop}{Block diagram of the aim loop: image (30 Hz, 0.1 s late)
-> fingertip -> base-frame ray (pose at the image stamp) -> One-Euro -> error in the
current camera frame -> P law + roll hold -> smoother -> DLS -> 125 Hz -> relay 500 Hz ->
velocity controller}{The finger-pointing camera aim loop.}
\SRC{admittance\_control/camera\_aim.py, scripts/camera\_aim\_node.py,
src/aim\_rate\_relay.cpp; thesis\_notes.md 2026-10-15/16.}

% -------------------------------------------------------------------------------------
\section{Perception and Registration}
\label{sec:mc_perception}

\subsection{Segmentation and Part Classification}
\label{sec:per_classify}
The clicked points (positive and negative) and the frozen colour image go to SAM2
\cite{ravi2024sam2}, which returns one mask per part; the masked depth becomes the part's
point cloud. Which CAD model it is follows from point-pair features
\cite{drost2010ppf,birdal2015ppf},
\begin{equation}
\mathbf{F}(\vect{m}_1,\vect{m}_2) = \big(\norm{\vect{d}},\ \angle(\vect{n}_1,\vect{d}),\
\angle(\vect{n}_2,\vect{d}),\ \angle(\vect{n}_1,\vect{n}_2)\big),\quad
\vect{d}=\vect{m}_2-\vect{m}_1,
\label{eq:ppf}
\end{equation}
voted against every model of the library; the best hypotheses are verified by how much
of the cloud they explain and how much of the model they cover.
\TODO{The verification score (coverage $\times$ explanation); 12 hypotheses; forcing the
model name on a re-registration.}
\SRC{notes/perception\_aciklama.md \S6--8; scripts\_in\_foundationpose/ppf\_classifier.py, fp\_server.py.}

\subsection{6D Pose Registration and Live Tracking}
\label{sec:per_register}
FoundationPose \cite{wen2024foundationpose} registers the CAD model to the masked RGB-D
frame (pose hypotheses, render-and-compare refinement, scoring) and then tracks it frame
to frame while the operator positions the part. Registration runs on the desktop GPU,
tracking on the laptop's; the camera's own motion is cancelled by looking up
Equation~\ref{eq:frame_chain} at each image's time stamp.
\TODO{Lost-track detection and re-registration locked to the tracked model's name; why
FoundationPose tracking replaced the ICP tracker (moving parts, CPU load).}
\SRC{notes/realtime\_fp.md; thesis\_notes.md 2026-10-02, 10-06.}
\figplaceholder{fig:tracking_split}{Desktop registers, laptop tracks: data flow, rates and
latencies}{The registration and tracking split between two GPUs.}

\subsection{Stationary Refinement and Saving}
\label{sec:per_stationary}
When a part comes to rest, its pose is averaged over the still frames (median
translation, chordal mean of the rotations) and refined by point-to-plane ICP against
the scene, minimising
\begin{equation}
E(\mathbf{R},\vect{t}) = \sum_i w_i\,\Big(\big(\mathbf{R}\vect{p}_i+\vect{t}-\vect{q}_i\big)^{\!\top}\vect{n}_i\Big)^2,
\qquad w_i = \exp\!\Big(-\frac{r_i^2}{2\nu^2}\Big),
\label{eq:p2plane}
\end{equation}
with Welsch weights \cite{zhang2022fricp,chen1992p2plane}. A pair counts only if the
normals agree,
\begin{equation}
(\mathbf{R}\,\vect{n}^{\mathrm{model}}_i)^{\top}\vect{n}^{\mathrm{scene}}_j \ge \cos 60^\circ ,
\label{eq:normal_gate}
\end{equation}
which stops a thin plate's model from straddling its own two faces. The saved part
becomes part of the static environment that the next part is registered against.
\TODO{Figure: the straddling plate before and after the gate (cross-section).}
\SRC{README Algorithms \S4--7; pose\_stats.py; icp\_pose\_refiner\_node.py.}

\subsection{Symmetric and Cut Parts}
\label{sec:per_symmetry}
A plain tube's rotation about its own axis is not observable, so ICP and the tracker each
return an arbitrary one. The correction is split into a swing (which tilts the axis) and
a twist (about it), and only the swing is kept:
\begin{equation}
\Rf{}{\mathrm{ref}}^{\top}\Rf{}{\mathrm{new}} = \mathbf{R}_{\mathrm{swing}}\,\mathbf{R}_{\mathrm{twist}}(\vect{a}),
\qquad \Rf{}{\mathrm{saved}} = \Rf{}{\mathrm{ref}}\,\mathbf{R}_{\mathrm{swing}} .
\label{eq:swing_twist}
\end{equation}
A tube cut at an angle (a mitre) is not symmetric, but its spin is still poorly
observed by the camera. It is fixed by physics instead: the part rests on its base, so
the spin is the one that makes the cut face sit flat on it,
\begin{equation}
\psi^{*} = \arg\min_{\psi}\ \Big(\max_{s} g(s;\psi) - \min_{s} g(s;\psi)\Big),
\label{eq:spin_settle}
\end{equation}
with $g(s;\psi)$ the gap between the cut face and the base along the contact.
\SRC{pose\_stats.remove\_twist; curved\_seams.settle\_cut\_spins; thesis\_notes.md 2026-10-15.}

% -------------------------------------------------------------------------------------
\section{Multi-View Close-Range Refinement}
\label{sec:mc_multiview}
A single view from the scan pose leaves millimetres of error where it matters most, at
the joints, which are seen at a distance and at grazing angles. With all parts saved, the
robot moves the camera to a few close views of the seams and all poses are refined
together.
\SRC{notes/multiview\_refine\_plan.md (design D1--D8); README \S15; multiview*.py.}

\subsection{View Planning}
\TODO{Targets on the seam regions; per candidate view the visibility of each target (range,
field of view, incidence, occlusion); greedy selection with diminishing returns; IK on
the branch, singularity and collision checks; the order of the views and the transits
between them.}
\figplaceholder{fig:views}{The planned views around an assembly (RViz) and the stacked
clouds coloured by view}{Planned close-range views.}

\subsection{Joint Refinement of Touching Parts}
For part $k$ with a correction $\boldsymbol{\xi}_k\in\mathfrak{se}(3)$ about its centre,
\begin{equation}
\min_{\boldsymbol{\xi}_k}\ \sum_{v}\sum_{i\in\mathcal{O}_k^v} \rho\big(r_{v,i}\big)^2
+ \boldsymbol{\xi}_k^{\top}\boldsymbol{\Sigma}^{-1}\boldsymbol{\xi}_k\,\sigma^2
+ w_o\sum_{c}\max\!\big(0,\,-\mathrm{sd}_c-\varepsilon\big)^2 ,
\label{eq:mv_objective}
\end{equation}
where $\mathcal{O}_k^v$ are the points of view $v$ that part $k$ owns, the second term is a
prior on the correction, and the third penalises overlap with the neighbouring parts,
$\mathrm{sd}_c$ being the signed distance of contact sample $c$ to its neighbour. The
neighbours are their true shapes, not bounding boxes: a tube's bore is hollow, a cut
tube ends at its cut. The parts take turns (block coordinate descent).
\TODO{Ownership with a dead band; the signed distance of each primitive (slab, tube with
its bore, cut tube, swept band); two rounds.}

\subsection{Observability}
With the information matrix $\mathbf{H}=\mathbf{J}^{\top}\mathbf{W}\mathbf{J}$ and its
eigenvalues $\lambda$, a direction is considered measured when
\begin{equation}
\sigma_{\mathrm{dir}} = \frac{\sigma_{\mathrm{noise}}}{\sqrt{\lambda}} \le 0.5\ \mathrm{mm}.
\label{eq:observability}
\end{equation}
Directions that are not measured keep their prior value.

\subsection{Online Estimation of the Camera's Extrinsic Translation}
\label{sec:mv_selfcal}
An error $\vect{d}$ in the extrinsic translation (camera frame) shifts view $v$'s cloud by
$\mathbf{R}_v\vect{d}$, differently in every view. Estimated jointly with the poses,
\begin{equation}
r_{v,i} = \vect{n}_i^{\top}\big(\mathbf{T}\vect{p}_i - \vect{q}_i + \mathbf{R}_v\vect{d}\big),
\qquad
\mathbf{S} = \mathbf{N}_{dd} - \mathbf{N}_{dx}\mathbf{N}_{xx}^{+}\mathbf{N}_{xd},
\label{eq:extrinsic_d}
\end{equation}
where the Schur complement $\mathbf{S}$ is the information about $\vect{d}$ left after
the poses are accounted for.
\TODO{Removing $\mathbf{R}_v\vect{d}$ from the views and $\mathbf{R}_{\mathrm{scan}}\vect{d}$
from the saved poses; why the joint fit alone does not remove $\vect{d}$.}

\subsection{Acceptance}
\label{sec:mv_accept}
A refinement is applied when it stays within the limits; the rules are consistent with
that principle:
\begin{table}[htbp]
\caption{When a multi-view correction is refused.}
\label{tab:mv_acceptance}
\centering\small
\begin{tabular}{@{}ll@{}}
\toprule
Check & Refused when \\
\midrule
Correction size & $>10$ mm or $>5^\circ$ (a pipe: the tilt of its axis, not its spin) \\
Overlap with neighbours & worse than the saved assembly's by more than 0.2 mm \\
Fit & \textit{to fill: fitness; an unmeasured relative change} \\
\bottomrule
\end{tabular}
\end{table}
\TODO{Complete the acceptance rules from multiview\_refine.judge.}

% -------------------------------------------------------------------------------------
\section{Seams from the Registered Assembly}
\label{sec:mc_seams}
Seams are not detected in the point cloud: they are computed from the registered CAD
models. Each part is represented as one of WeldSet's exact primitives
(Chapter~\ref{chp:method_weldset}), and the seam rule of WeldSet
(\S\ref{sec:ws_d4}) is applied to the primitives at their registered poses. The cell's
seams are thus the same exact geometry that WeldSet uses as ground truth, evaluated at
the measured poses.

\subsection{CAD Parts as Primitives}
\label{sec:mc_registry}
A registry maps every CAD model to a primitive: a slab (plate), a tube (with its bore,
and its cuts), or a swept slab (a band along a spine curve, fitted as a B-spline). A fit
is accepted when it is verified both ways, the CAD surface against the primitive and the
primitive against the CAD surface, within a fraction of a millimetre, and the volumes
agree; a registry entry whose CAD file has changed since the fit is detected and refused.
\SRC{weldgen\_registry.py (derive\_box, derive\_tube, derive\_extrusion, verify\_both\_ways);
part\_shapes.stale\_entries.}

\subsection{Planar Seams}
Two planar faces with unit normals $\vect{n}_1,\vect{n}_2$ through points
$\vect{p}_1,\vect{p}_2$ meet on the line
\begin{equation}
\vect{d} = \frac{\vect{n}_1\times\vect{n}_2}{\norm{\vect{n}_1\times\vect{n}_2}},\qquad
\vect{x}(s) = \vect{x}_0 + s\,\vect{d},\quad
\vect{n}_i^{\top}\vect{x}_0 = \vect{n}_i^{\top}\vect{p}_i\ (i=1,2),
\label{eq:plane_intersection}
\end{equation}
clipped to where both faces exist. Because the parts are registered, not in perfect
contact, the rule uses a pose tolerance instead of a contact tolerance.

\subsection{Curved Seams}
\label{sec:mc_curved}
A tube standing on a plate meets it along a plane--cylinder intersection, an ellipse.
With the tube axis $\vect{a}$ through $\vect{c}$, radius $r$, orthonormal $\vect{u},\vect{v}\perp\vect{a}$,
and the plate's top plane $\vect{n}^{\top}\vect{x}=\vect{n}^{\top}\vect{p}_0$,
\begin{equation}
\begin{aligned}
\vect{x}(\theta) &= \vect{c} + r(\cos\theta\,\vect{u}+\sin\theta\,\vect{v}) + h(\theta)\,\vect{a},\\
h(\theta) &= \frac{\vect{n}^{\top}(\vect{p}_0-\vect{c}) - r\,(\vect{n}^{\top}\vect{u}\cos\theta+\vect{n}^{\top}\vect{v}\sin\theta)}{\vect{n}^{\top}\vect{a}} .
\end{aligned}
\label{eq:ellipse_seam}
\end{equation}
A branch tube on a main tube meets it along a saddle (cylinder--cylinder intersection);
a curved band on a plate meets it along the offsets of its spine.
\TODO{The saddle in general position (the branch at an angle); the band's two toe
curves; a member must not lie in the face's plane (minimum angle); the frame per point
(normals of both members, approach) and why it rotates along a curved seam.}
\SRC{curved\_seams.py (\_tube\_on\_slab, \_tube\_on\_tube, \_band\_on\_slab);
weldgen/curves.py (ellipse\_from\_plane\_cylinder, saddle\_from\_cylinders).}
\figplaceholder{fig:mc_seam_kinds}{The three curved cases on the bench parts: pipe on
plate (ellipse), branch on pipe (saddle), band on plate (spine offsets), with the
per-point frames}{Curved seams computed from registered primitives.}

\subsection{Fit-Up Gap and Weldability}
\label{sec:mc_fitup}
Between registered parts the gap along a seam is measured, and a seam whose gap exceeds
what a welder could bridge is reported, not silently dropped. The admissible root gap
follows ISO 5817 \cite{iso5817}:
\begin{equation}
h_{\max} = c_0 + c_1\,a,\quad a = 0.7\,t:\ \ (c_0,c_1) = (0.5,0.1)_{B},\ (0.5,0.2)_{C},\ (1,0.3)_{D}\ \mathrm{mm}.
\label{eq:iso_gap}
\end{equation}
\TODO{Verify Eq.~\ref{eq:iso_gap} and its caps against Table~1, no.~617 of the standard.
Near misses (up to twice the tolerance) reported with the reason.}

% -------------------------------------------------------------------------------------
\section{Tack Placement and Ordering}
\label{sec:mc_tacks}
Tacks follow WeldSet's tack rule (\S\ref{sec:ws_tacks}, Equations~\ref{eq:tackrule_len}
and~\ref{eq:tackrule_n}) applied to the computed seams, open or closed: both ends tacked,
the rest evenly spaced. The rule also fixes the order in which they would be welded: the
ends first, then by bisection, alternating between seams so that heat does not
accumulate in one place.
\TODO{OPEN: tack ordering by a rough heat simulation (todo.md, 2026-10-16): each tack a
heat source, superposed cooling kernels, the next tack the coolest; compared with the
rule's order. In this thesis or future work?}
\TODO{OPEN DECISION (THESIS\_PLAN.md \#1): robot-aware placement (a quality field over arc
length $\times$ roll and an exact DP for $k$ tacks). Note: the per-tack roll DP of
\S\ref{sec:mc_reach} already chooses the robot's side of it.}

% -------------------------------------------------------------------------------------
\section{Motion Planning and Execution}
\label{sec:mc_motion}

\subsection{Collision Model}
\label{sec:mc_collision}
The arm is a chain of capsules on its joint frames; the pen tool is its CAD envelope
(capsules for the pen and holder, boxes for the camera body); the parts are their
primitives as boxes (a tube as rotated boxes round its axis, a cut tube as sector boxes,
a band as a chain of boxes); the table is the touched plane, tilted; fixtures are boxes
entered by pen-tip points. Distances are exact between capsules and between a capsule and
a box; a box against a box is an intersection test only.
\TODO{Clearance on tacks vs.\ transits; which contacts are allowed (the pen tip with the
part at the tack); the C++ port of the model, proven equal to the Python one.}
\SRC{collision.py docstring; src/transit\_cpp.cpp; test\_transit\_cpp.py.}
\figplaceholder{fig:mc_collision}{The collision model around a tube-on-plate assembly:
arm capsules, tool envelope, part boxes, table plane}{The collision model.}

\subsection{Reachability}
\label{sec:mc_reach}
Inverse kinematics is solved on one branch, the elbow-up configuration of the home pose
(the signs of the lift, elbow and second wrist joints are locked). On a straight seam one
pen roll and tilt is chosen per seam, maximising the smallest clearance over the
approach pose, the tack pose and the straight descent between them. On a curved seam the
pen's frame rotates with the seam, so a roll is chosen per tack, jointly over the seam by
dynamic programming (a Viterbi recursion \cite{forney1973viterbi}):
\begin{equation}
\min_{\varphi_1,\dots,\varphi_n}\ \sum_{k=1}^{n}\Big[
\norm{\mathbf{W}\,\big(\vect{q}_k(\varphi_k)-\vect{q}_{k-1}(\varphi_{k-1})\big)}
+ \kappa\,\max\!\big(0,\ \delta_0-\delta_k(\varphi_k)\big)\Big],
\label{eq:roll_dp}
\end{equation}
with $\delta_k$ the clearance of tack $k$ at roll $\varphi_k$, $\delta_0=10$~mm and
$\kappa=0.15$~rad/mm. Every pose is taken as the $2\pi$-equivalent that keeps the last
wrist joint within $\pm180^\circ$ of home: the camera's cable wraps round the wrist, and
plain joint continuity wound it far enough for its pull to read as a contact.
\SRC{tack\_reach.py (plan\_tacks, \_plan\_per\_tack, to\_cable\_window); thesis\_notes.md 2026-10-16.}
\figplaceholder{fig:mc_roll_dp}{The trellis of rolls along a curved seam: feasible rolls
per tack, the chosen path}{Choosing a pen roll per tack on a curved seam.}

\subsection{Transit Planning}
\label{sec:mc_transit}
Transits start and end at departure points, 100~mm out along each tack's pen axis, where
the tool is in open space. A transit is tried, in order, as: the straight joint-space
edge; the tool's straight Cartesian line (orientation interpolated by slerp), then the
same line bowed outward as a quadratic B\'ezier,
\begin{equation}
\vect{x}(s) = (1-s)^2\,\vect{x}_0 + 2s(1-s)\,(\vect{x}_m + \beta\,\vect{b}) + s^2\,\vect{x}_1,
\qquad \beta\in\{0, 30, 60, 100\}\ \mathrm{mm};
\label{eq:cartesian_bow}
\end{equation}
and only then a sampling-based planner: OMPL's AnytimePathShortening
\cite{luo2014aps,sucan2012ompl} (parallel RRT-Connect \cite{kuffner2000rrtconnect},
hybridisation and shortcutting within a time budget), in scaled joint coordinates
$y_j=w_j q_j$ so that path length is weighted joint travel,
\begin{equation}
c(\text{path}) = \sum_{k}\norm{\mathbf{W}(\vect{q}_{k+1}-\vect{q}_k)}_2,\qquad
\mathbf{W}=\operatorname{diag}(2,\,2,\,1.5,\,1,\,1,\,0.5).
\label{eq:transit_cost}
\end{equation}
The path is smoothed, re-checked against the collision model and timed with
time-optimal trajectory generation \cite{kunz2012totg}.
\TODO{Why departure points (planning from beside a plate); why the Cartesian line first
(round a pipe it is short and the same every time); the planning margin capped by the
end poses' own room.}
\SRC{marking.py (transit\_path, cartesian\_transit, aps\_path, departure\_chain);
notes/aps\_transit\_plan.md.}

\subsection{Force-Gated Descent and Marking}
\label{sec:mc_marking}
From the departure point the pen moves along its own axis: fast to the approach point
(35~mm before the tack), slower to the tack, slowest in the last millimetres, and on past
the registered point by an overshoot, in case the real surface is farther. The descent
stops at the first sample where the force, with its bias measured at rest just before,
exceeds the touch threshold:
\begin{equation}
\norm{\vect{F}(t)-\vect{F}_{\mathrm{bias}}} > F_{\mathrm{touch}} = 1.5\ \mathrm{N},
\qquad
\delta_c = (\vect{p}_{\mathrm{reg}}-\vect{p}_{\mathrm{tip}})^{\top}\vect{a},
\label{eq:touch}
\end{equation}
where $\delta_c$, the contact depth, says how far before ($+$) or past ($-$) the registered
point the surface was met: an independent measurement of the registration along the pen
axis. The overshoot ends at the table, since no surface lies beyond it. The mark is a
dot, or a stroke along the seam with the depth corrected from the force between chunks.
\TODO{The speed profile and why (an unsprung pen on a rigid part: force grows with
speed); the retract; the 8 N abort; strokes on curved seams.}
\SRC{tack\_marking\_node.py, marking.py (descent\_check, time\_descent), motion.py.}
\figplaceholder{fig:marking_sequence}{One tack: approach, descent, contact, stroke,
retract; joint trajectory and the force along the pen}{One tack, from transit to retract.}

% -------------------------------------------------------------------------------------
\section{Implementation Details}
\label{sec:mc_impl}

\subsection{ROS 2 Architecture}
\TODO{The nodes and what feeds what: camera, bridge (to the pose server), ICP / assembly
node (tracking, saving, seams, refinement), gesture and aim nodes, reachability, marking;
services as the steps of the pipeline; the static environment cloud.}
\figplaceholder{fig:architecture}{Nodes, topics and services; the desktop (registration
server in Docker) and the laptop (everything else)}{Software architecture of the cell.}
\SRC{README ``How the pieces connect'' and the node table.}

\subsection{Performance-Critical Components}
\TODO{What runs in C++ and why: the collision model and the planners (one plan from minutes
to seconds), the time-parameterisation service, the 500 Hz command relay.}
